%% relecture.tex : la version de relecture d'un document de chantier (make relecture), pour lire
%% sur tablette sans aller-retour vers la bibliographie. Les pages restent celles du document, au
%% trait près ; sous chaque page qui cite, ou qui a des tournures repérées, la page est allongée
%% d'une bande grise qui donne :
%%   - les références citées sur cette page : numéro (lien vers la bibliographie), auteurs, titre,
%%     lieu, année ;
%%   - les tournures à revoir, numérotées comme dans la marge de la page, avec leur raison.
%% Dans la page, une tournure repérée est surlignée et soulignée, et son numéro est dans la marge :
%% fond jaune et trait plein pour une tournure d'assistant, fond bleu clair et trait plein pour une
%% faute ou un calque, trait pointillé gris sans fond pour un indice faible (« à vérifier »). La
%% couleur n'est jamais seule : le trait et la bande disent aussi de quoi il s'agit.
%%
%% Chargé avant \documentclass par SCRIPTS/relecture.sh, qui fait les passages :
%%   pdflatex -jobname=<job> "\AddToHook{begindocument/before}{%% relecture.tex : la version de relecture d'un document de chantier (make relecture), pour lire
%% sur tablette sans aller-retour vers la bibliographie. Les pages restent celles du document, au
%% trait près ; sous chaque page qui cite, ou qui a des tournures repérées, la page est allongée
%% d'une bande grise qui donne :
%%   - les références citées sur cette page : numéro (lien vers la bibliographie), auteurs, titre,
%%     lieu, année ;
%%   - les tournures à revoir, numérotées comme dans la marge de la page, avec leur raison.
%% Dans la page, une tournure repérée est surlignée et soulignée, et son numéro est dans la marge :
%% fond jaune et trait plein pour une tournure d'assistant, fond bleu clair et trait plein pour une
%% faute ou un calque, trait pointillé gris sans fond pour un indice faible (« à vérifier »). La
%% couleur n'est jamais seule : le trait et la bande disent aussi de quoi il s'agit.
%%
%% Chargé avant \documentclass par SCRIPTS/relecture.sh, qui fait les passages :
%%   pdflatex -jobname=<job> "\AddToHook{begindocument/before}{%% relecture.tex : la version de relecture d'un document de chantier (make relecture), pour lire
%% sur tablette sans aller-retour vers la bibliographie. Les pages restent celles du document, au
%% trait près ; sous chaque page qui cite, ou qui a des tournures repérées, la page est allongée
%% d'une bande grise qui donne :
%%   - les références citées sur cette page : numéro (lien vers la bibliographie), auteurs, titre,
%%     lieu, année ;
%%   - les tournures à revoir, numérotées comme dans la marge de la page, avec leur raison.
%% Dans la page, une tournure repérée est surlignée et soulignée, et son numéro est dans la marge :
%% fond jaune et trait plein pour une tournure d'assistant, fond bleu clair et trait plein pour une
%% faute ou un calque, trait pointillé gris sans fond pour un indice faible (« à vérifier »). La
%% couleur n'est jamais seule : le trait et la bande disent aussi de quoi il s'agit.
%%
%% Chargé avant \documentclass par SCRIPTS/relecture.sh, qui fait les passages :
%%   pdflatex -jobname=<job> "\AddToHook{begindocument/before}{%% relecture.tex : la version de relecture d'un document de chantier (make relecture), pour lire
%% sur tablette sans aller-retour vers la bibliographie. Les pages restent celles du document, au
%% trait près ; sous chaque page qui cite, ou qui a des tournures repérées, la page est allongée
%% d'une bande grise qui donne :
%%   - les références citées sur cette page : numéro (lien vers la bibliographie), auteurs, titre,
%%     lieu, année ;
%%   - les tournures à revoir, numérotées comme dans la marge de la page, avec leur raison.
%% Dans la page, une tournure repérée est surlignée et soulignée, et son numéro est dans la marge :
%% fond jaune et trait plein pour une tournure d'assistant, fond bleu clair et trait plein pour une
%% faute ou un calque, trait pointillé gris sans fond pour un indice faible (« à vérifier »). La
%% couleur n'est jamais seule : le trait et la bande disent aussi de quoi il s'agit.
%%
%% Chargé avant \documentclass par SCRIPTS/relecture.sh, qui fait les passages :
%%   pdflatex -jobname=<job> "\AddToHook{begindocument/before}{\input{relecture}}\input{<source>}"
%% Références : chaque citation écrit \rlcite{clé}{page} dans le .aux, à la sortie de la page, donc
%% lu au passage suivant ; les entrées courtes viennent de <job>-refs.bbl, que bibtex produit avec
%% relecture.bst sur les clés citées. Marche avec natbib (acmart, LIPIcs) et avec le \cite de
%% LaTeX, hyperref chargé ou non ; pas avec biblatex.
%% Tournures : au dernier passage seulement (\rlavectournures défini), <job>-tournures.aux, écrit
%% par SCRIPTS/relecture-tournures.py d'après le texte du PDF (pdftotext -bbox) : des rectangles en
%% points PostScript (bp) depuis le coin haut gauche de la page, et les lignes de la bande. Pour
%% lui, le .aux donne la hauteur de la page sans bande (\rlhauteur) et les pages de la
%% bibliographie (\rlbiblio), où rien n'est cherché.
%% Ne part jamais vers le finalisé : le PDF de relecture est <document>-relecture.pdf.

\providecommand\rltitre{R\'ef\'erences cit\'ees sur cette page}
\providecommand\rltitretournures{Tournures \`a revoir}
\providecommand\rlhauteur[1]{}
\providecommand\rlbiblio[2]{}

\RequirePackage{xcolor}
\definecolor{rlfort}{RGB}{255,232,120}   % fond : tournure d'assistant (jaune)
\definecolor{rlfaute}{RGB}{196,222,255}  % fond : faute ou calque (bleu clair)
\definecolor{rltrait}{RGB}{25,55,145}    % traits pleins et numéros (bleu foncé)
\definecolor{rlgris}{gray}{0.45}         % indice faible : pointillé, numéro et ligne de la bande

\ExplSyntaxOn
\makeatletter

\prop_new:N \g__rl_entree_prop     % clé -> {Auteurs.}{Titre.}{Lieu, année.}
\prop_new:N \g__rl_num_prop        % clé -> étiquette dans la bibliographie
\prop_new:N \g__rl_rang_prop       % clé -> rang dans la bibliographie
\int_new:N \g__rl_rang_int
\box_new:N \l__rl_etiquette_box
\prop_new:N \g__rl_page_prop       % page absolue -> clés citées, séparées par des virgules
\prop_new:N \g__rl_dessin_prop     % page absolue -> fonds, traits et numéros des tournures
\prop_new:N \g__rl_tournure_prop   % page absolue -> lignes « Tournures à revoir » de la bande
\box_new:N \g__rl_bande_box
\tl_new:N \g__rl_dessin_tl         % dessin de la page en cours de sortie
\dim_new:N \g__rl_bande_dim        % hauteur ajoutée à la page courante
\dim_new:N \g__rl_hauteur_dim      % hauteur de la page sans bande
\seq_new:N \l__rl_cles_seq
\tl_new:N \l__rl_tmp_tl
\tl_new:N \l__rl_refs_tl
\tl_new:N \l__rl_tour_tl
\tl_new:N \l__rl_page_tl
\dim_const:Nn \c__rl_marge_dim { 14pt }
\cs_generate_variant:Nn \prop_get:NnNF { Ne }
\cs_generate_variant:Nn \prop_get:NnNT { Ne }

%% Ajoute du code à la valeur d'une clé (une page) d'une propriété globale.
\cs_new_protected:Npn \__rl_ajouter:Nnn #1#2#3
  {
    \prop_get:NnNF #1 {#2} \l__rl_tmp_tl { \tl_clear:N \l__rl_tmp_tl }
    \tl_put_right:Nn \l__rl_tmp_tl {#3}
    \prop_gput:NnV #1 {#2} \l__rl_tmp_tl
  }

%% Les entrées courtes, lues dans <job>-refs.bbl.
\cs_new_protected:Npn \rlentry #1#2#3#4
  { \prop_gput:Nnn \g__rl_entree_prop {#1} { {#2} {#3} {#4} } }

%% \rlcite{clé}{page}, lu dans le .aux : la clé s'ajoute aux clés de la page.
\cs_new_protected:Npn \rlcite #1#2
  {
    \prop_get:NnNF \g__rl_page_prop {#2} \l__rl_tmp_tl { \tl_clear:N \l__rl_tmp_tl }
    \clist_if_in:NnF \l__rl_tmp_tl {#1}
      {
        \clist_put_right:Nn \l__rl_tmp_tl {#1}
        \prop_gput:NnV \g__rl_page_prop {#2} \l__rl_tmp_tl
      }
  }

%% Les tournures, lues dans <job>-tournures.aux (coordonnées en bp, l'origine en haut à gauche,
%% y vers le haut : négatif dans la page).
%%   \rlfond{page}{couleur}{x}{y}{largeur}{hauteur}   fond sous les mots, (x, y) coin bas gauche
%%   \rltrait{page}{x}{y}{largeur}                    trait plein sous les mots
%%   \rlpointille{page}{x}{y}{nombre}                 trait pointillé gris, un pas tous les 3 bp
%%   \rlnum{page}{couleur}{x}{y}{numéros}             numéros dans la marge, alignés à droite en x
%%   \rltournure{page}{n}{genre}{texte}{raison}       ligne de la bande ; genre f, c ou v
\cs_new_protected:Npn \rlfond #1#2#3#4#5#6
  {
    \__rl_ajouter:Nnn \g__rl_dessin_prop {#1}
      { \put (#3,#4) { \textcolor {#2} { \rule {#5\unitlength} {#6\unitlength} } } }
  }
\cs_new_protected:Npn \rltrait #1#2#3#4
  {
    \__rl_ajouter:Nnn \g__rl_dessin_prop {#1}
      { \put (#2,#3) { \textcolor {rltrait} { \rule {#4\unitlength} {0.6\unitlength} } } }
  }
\cs_new_protected:Npn \rlpointille #1#2#3#4
  {
    \__rl_ajouter:Nnn \g__rl_dessin_prop {#1}
      { \multiput (#2,#3) (3,0) {#4} { \textcolor {rlgris} { \rule {1.6\unitlength} {0.5\unitlength} } } }
  }
\cs_new_protected:Npn \rlnum #1#2#3#4#5
  {
    \__rl_ajouter:Nnn \g__rl_dessin_prop {#1}
      { \put (#3,#4) { \makebox (0,0) [r] { \normalfont \sffamily \bfseries \tiny \textcolor {#2} {#5} } } }
  }
\cs_new_protected:Npn \rltournure #1#2#3#4#5
  { \__rl_ajouter:Nnn \g__rl_tournure_prop {#1} { \__rl_ligne_tournure:nnnn {#2} {#3} {#4} {#5} } }

%% Le numéro de chaque clé, pris au passage de \bibcite dans le .aux : natbib écrit
%% {numéro}{année}{{auteurs}}{{...}}, LaTeX l'étiquette seule (un numéro, ou ABK+06 en style
%% alpha), beamer l'étiquette dans un \hyperlink vers sa page de bibliographie. Les \bibcite
%% viennent dans l'ordre de la bibliographie : leur rang ordonne les lignes de la bande.
\cs_new_eq:NN \__rl_bibcite_orig:nn \bibcite
\cs_gset_protected:Npn \bibcite #1#2
  {
    \__rl_bibcite_orig:nn {#1} {#2}
    \prop_if_in:NnF \g__rl_rang_prop {#1}
      {
        \int_gincr:N \g__rl_rang_int
        \prop_gput:Nne \g__rl_rang_prop {#1} { \int_use:N \g__rl_rang_int }
      }
    \cs_if_exist:NTF \NAT@parse
      { \prop_gput:Nno \g__rl_num_prop {#1} { \tl_head:n {#2} } }
      { \prop_gput:Nnn \g__rl_num_prop {#1} {#2} }
  }

%% Écrit #1 dans le .aux à la sortie de la page, avec la page absolue en dernier argument.
\cs_new_protected:Npn \__rl_ecrire_page:n #1
  {
    \protected@write \@auxout {}
      { #1 { \exp_not:N \the \exp_not:N \ReadonlyShipoutCounter } }
  }

%% Chaque clé citée écrit \rlcite{clé}{page absolue} à la sortie de sa page.
\cs_new_protected:Npn \__rl_note:n #1
  {
    \@bsphack
    \clist_map_inline:nn {#1}
      { \__rl_ecrire_page:n { \token_to_str:N \rlcite {##1} } }
    \@esphack
  }

%% Branchement, une fois tous les paquets en place (natbib choisit son \@citex au début du
%% document) : toutes les citations passent par \@citex, avec deux arguments optionnels sous
%% natbib, un seul avec le \cite de LaTeX ; la liste des clés est le dernier argument. Les clés
%% sont notées même indéfinies, dès le premier passage. (Dans une macro : le code d'un crochet ne
%% dédouble pas les #.) Le .aux reçoit aussi la hauteur de la page, pour relecture-tournures.py.
\cs_new_protected:Npn \__rl_brancher:
  {
    \cs_gset_eq:NN \__rl_citex_orig:w \@citex
    \cs_if_exist:NTF \NAT@parse
      {
        \cs_gset_protected:Npn \@citex [##1] [##2] ##3
          { \__rl_citex_orig:w [{##1}] [{##2}] {##3} \__rl_note:n {##3} }
      }
      {
        \cs_gset_protected:Npn \@citex [##1] ##2
          { \__rl_citex_orig:w [{##1}] {##2} \__rl_note:n {##2} }
      }
    \iow_now:Ne \@auxout { \token_to_str:N \rlhauteur { \dim_use:N \paperheight } }
  }
\AddToHook { begindocument / end } { \__rl_brancher: }

%% Les pages de la bibliographie, où relecture-tournures.py ne cherche rien : la page de son
%% premier paragraphe (le titre ou la première entrée) et celle de sa fin. Rien n'est écrit entre
%% deux paragraphes, où une note dans la liste verticale changerait les espaces.
\AddToHook { env / thebibliography / begin }
  { \AddToHookNext { para / begin } { \__rl_ecrire_page:n { \token_to_str:N \rlbiblio {debut} } } }
\AddToHook { env / thebibliography / end }
  { \@bsphack \__rl_ecrire_page:n { \token_to_str:N \rlbiblio {fin} } \@esphack }

%% La bande d'une page : les références, dans l'ordre de la bibliographie, puis les tournures.
\cs_new:Npn \__rl_num:n #1
  { \prop_if_in:NnTF \g__rl_rang_prop {#1} { \prop_item:Nn \g__rl_rang_prop {#1} } { 0 } }

%% L'étiquette entre crochets, lien vers la bibliographie ; celle de beamer porte déjà le sien
%% (deux liens emboîtés arrêteraient pdfTeX).
\cs_new_protected:Npn \__rl_etiquette:n #1
  {
    \prop_get:NnNF \g__rl_num_prop {#1} \l__rl_tmp_tl { \tl_set:Nn \l__rl_tmp_tl { ? } }
    \cs_if_exist:NTF \hyperlink
      {
        \tl_if_in:NnTF \l__rl_tmp_tl { \hyperlink }
          { [ \l__rl_tmp_tl ] }
          { \hyperlink { cite.#1 } { [ \l__rl_tmp_tl ] } }
      }
      { [ \l__rl_tmp_tl ] }
  }

\cs_new_protected:Npn \__rl_ligne:n #1
  {
    \hangindent = 2.9em \hangafter = 1 \noindent
    \hbox_set:Nn \l__rl_etiquette_box { \__rl_etiquette:n {#1} }
    \hbox_to_wd:nn { \dim_max:nn { 2.9em } { \box_wd:N \l__rl_etiquette_box + 0.5em } }
      { \box_use_drop:N \l__rl_etiquette_box \hfil }
    \prop_get:NnNTF \g__rl_entree_prop {#1} \l__rl_tmp_tl
      { \exp_after:wN \__rl_entree:nnn \l__rl_tmp_tl }
      { \texttt { \tl_to_str:n {#1} } }
    \par
  }
\cs_new_protected:Npn \__rl_entree:nnn #1#2#3
  { \tl_if_empty:nF {#1} { #1 \c_space_tl } \textit {#2} \tl_if_empty:nF {#3} { \c_space_tl #3 } }

%% Une ligne des tournures : numéro (celui de la marge), texte repéré, raison ; en gris pour un
%% indice faible.
\cs_new_protected:Npn \__rl_ligne_tournure:nnnn #1#2#3#4
  {
    \group_begin:
      \str_if_eq:nnT {#2} {v} { \color {rlgris} }
      \hangindent = 2.9em \hangafter = 1 \noindent
      \hbox_to_wd:nn { 2.9em }
        { \sffamily \bfseries \str_if_eq:nnF {#2} {v} { \color {rltrait} } #1 \hfil }
      \textit {#3} \c_space_tl : \c_space_tl #4
      \par
    \group_end:
  }

\cs_new_protected:Npn \__rl_bande:
  {
    \tl_if_empty:NF \l__rl_refs_tl
      {
        \exp_args:NNV \seq_set_from_clist:Nn \l__rl_cles_seq \l__rl_refs_tl
        \seq_sort:Nn \l__rl_cles_seq
          {
            \int_compare:nNnTF { \__rl_num:n {##1} } > { \__rl_num:n {##2} }
              { \sort_return_swapped: } { \sort_return_same: }
          }
      }
    \vbox_gset:Nn \g__rl_bande_box
      {
        \hsize = \textwidth \linewidth = \textwidth
        \normalfont \footnotesize \raggedright
        \parindent = 0pt \parskip = 1.5pt \relax
        \cs_set:Npn \\ { \relax \ignorespaces }
        \tl_if_empty:NF \l__rl_refs_tl
          {
            \noindent \textbf { \rltitre } \par
            \seq_map_function:NN \l__rl_cles_seq \__rl_ligne:n
          }
        \tl_if_empty:NF \l__rl_tour_tl
          {
            \tl_if_empty:NF \l__rl_refs_tl { \skip_vertical:n { 4pt } }
            \noindent \textbf { \rltitretournures } \par
            \l__rl_tour_tl
          }
      }
  }

%% Avant la sortie de chaque page (le compteur vaut alors le numéro de la page précédente) :
%% la bande, le dessin des tournures, et la hauteur de la page.
\AddToHook { shipout / before }
  {
    \dim_compare:nNnT { \g__rl_hauteur_dim } = { 0pt }
      { \dim_gset:Nn \g__rl_hauteur_dim { \paperheight } }
    \box_gclear:N \g__rl_bande_box
    \dim_gzero:N \g__rl_bande_dim
    \tl_set:Ne \l__rl_page_tl { \int_eval:n { \ReadonlyShipoutCounter + 1 } }
    \prop_get:NeNF \g__rl_page_prop { \l__rl_page_tl } \l__rl_refs_tl { \tl_clear:N \l__rl_refs_tl }
    \prop_get:NeNF \g__rl_tournure_prop { \l__rl_page_tl } \l__rl_tour_tl { \tl_clear:N \l__rl_tour_tl }
    \prop_get:NeNF \g__rl_dessin_prop { \l__rl_page_tl } \l__rl_tmp_tl { \tl_clear:N \l__rl_tmp_tl }
    \tl_gset_eq:NN \g__rl_dessin_tl \l__rl_tmp_tl
    \bool_lazy_and:nnF { \tl_if_empty_p:N \l__rl_refs_tl } { \tl_if_empty_p:N \l__rl_tour_tl }
      {
        \__rl_bande:
        \dim_gset:Nn \g__rl_bande_dim
          { \box_ht_plus_dp:N \g__rl_bande_box + 2 \c__rl_marge_dim }
      }
    \cs_if_exist:NT \pdfpageheight
      { \global \pdfpageheight = \dim_eval:n { \g__rl_hauteur_dim + \g__rl_bande_dim } \relax }
  }

%% Sous le texte de la page : les fonds et les traits des tournures, puis la bande, sur fond gris
%% clair et sur toute la largeur, son texte aligné sur la colonne.
\cs_new:Npn \__rl_pt:n #1 { \dim_to_decimal_in_unit:nn {#1} { 1pt } }
\AddToHook { shipout / background }
  {
    \tl_if_empty:NF \g__rl_dessin_tl
      {
        \group_begin:
          \setlength \unitlength { 1bp }
          \g__rl_dessin_tl
        \group_end:
      }
    \box_if_empty:NF \g__rl_bande_box
      {
        \setlength \unitlength { 1pt }
        \put ( 0 , \__rl_pt:n { - \g__rl_hauteur_dim - \g__rl_bande_dim } )
          { \textcolor { black!6 } { \rule { \paperwidth } { \g__rl_bande_dim } } }
        \put ( 0 , \__rl_pt:n { - \g__rl_hauteur_dim } )
          { \textcolor { black!35 } { \rule { \paperwidth } { 0.6pt } } }
        \put
          (
            \__rl_pt:n
              {
                1in + \hoffset +
                \bool_lazy_and:nnTF { \legacy_if_p:n { @twoside } } { \int_if_even_p:n { \c@page } }
                  { \evensidemargin } { \oddsidemargin }
              } ,
            \__rl_pt:n
              { - \g__rl_hauteur_dim - \c__rl_marge_dim - \box_ht:N \g__rl_bande_box }
          )
          { \box_use:N \g__rl_bande_box }
      }
  }

\makeatother
\ExplSyntaxOff

%% Les entrées courtes (absentes au premier passage) et, au dernier passage, les tournures, lues
%% hors de la syntaxe expl3 qui ignore les espaces. Comme pour la bibliographie, & et \\ sont
%% neutralisés (un champ journal en contient parfois, dans la biblio de référence).
\begingroup
  \catcode`\&=12
  \InputIfFileExists{\jobname-refs.bbl}{}{}
\endgroup
\ifdefined\rlavectournures
  \InputIfFileExists{\jobname-tournures.aux}{}{}
\fi
}\input{<source>}"
%% Références : chaque citation écrit \rlcite{clé}{page} dans le .aux, à la sortie de la page, donc
%% lu au passage suivant ; les entrées courtes viennent de <job>-refs.bbl, que bibtex produit avec
%% relecture.bst sur les clés citées. Marche avec natbib (acmart, LIPIcs) et avec le \cite de
%% LaTeX, hyperref chargé ou non ; pas avec biblatex.
%% Tournures : au dernier passage seulement (\rlavectournures défini), <job>-tournures.aux, écrit
%% par SCRIPTS/relecture-tournures.py d'après le texte du PDF (pdftotext -bbox) : des rectangles en
%% points PostScript (bp) depuis le coin haut gauche de la page, et les lignes de la bande. Pour
%% lui, le .aux donne la hauteur de la page sans bande (\rlhauteur) et les pages de la
%% bibliographie (\rlbiblio), où rien n'est cherché.
%% Ne part jamais vers le finalisé : le PDF de relecture est <document>-relecture.pdf.

\providecommand\rltitre{R\'ef\'erences cit\'ees sur cette page}
\providecommand\rltitretournures{Tournures \`a revoir}
\providecommand\rlhauteur[1]{}
\providecommand\rlbiblio[2]{}

\RequirePackage{xcolor}
\definecolor{rlfort}{RGB}{255,232,120}   % fond : tournure d'assistant (jaune)
\definecolor{rlfaute}{RGB}{196,222,255}  % fond : faute ou calque (bleu clair)
\definecolor{rltrait}{RGB}{25,55,145}    % traits pleins et numéros (bleu foncé)
\definecolor{rlgris}{gray}{0.45}         % indice faible : pointillé, numéro et ligne de la bande

\ExplSyntaxOn
\makeatletter

\prop_new:N \g__rl_entree_prop     % clé -> {Auteurs.}{Titre.}{Lieu, année.}
\prop_new:N \g__rl_num_prop        % clé -> étiquette dans la bibliographie
\prop_new:N \g__rl_rang_prop       % clé -> rang dans la bibliographie
\int_new:N \g__rl_rang_int
\box_new:N \l__rl_etiquette_box
\prop_new:N \g__rl_page_prop       % page absolue -> clés citées, séparées par des virgules
\prop_new:N \g__rl_dessin_prop     % page absolue -> fonds, traits et numéros des tournures
\prop_new:N \g__rl_tournure_prop   % page absolue -> lignes « Tournures à revoir » de la bande
\box_new:N \g__rl_bande_box
\tl_new:N \g__rl_dessin_tl         % dessin de la page en cours de sortie
\dim_new:N \g__rl_bande_dim        % hauteur ajoutée à la page courante
\dim_new:N \g__rl_hauteur_dim      % hauteur de la page sans bande
\seq_new:N \l__rl_cles_seq
\tl_new:N \l__rl_tmp_tl
\tl_new:N \l__rl_refs_tl
\tl_new:N \l__rl_tour_tl
\tl_new:N \l__rl_page_tl
\dim_const:Nn \c__rl_marge_dim { 14pt }
\cs_generate_variant:Nn \prop_get:NnNF { Ne }
\cs_generate_variant:Nn \prop_get:NnNT { Ne }

%% Ajoute du code à la valeur d'une clé (une page) d'une propriété globale.
\cs_new_protected:Npn \__rl_ajouter:Nnn #1#2#3
  {
    \prop_get:NnNF #1 {#2} \l__rl_tmp_tl { \tl_clear:N \l__rl_tmp_tl }
    \tl_put_right:Nn \l__rl_tmp_tl {#3}
    \prop_gput:NnV #1 {#2} \l__rl_tmp_tl
  }

%% Les entrées courtes, lues dans <job>-refs.bbl.
\cs_new_protected:Npn \rlentry #1#2#3#4
  { \prop_gput:Nnn \g__rl_entree_prop {#1} { {#2} {#3} {#4} } }

%% \rlcite{clé}{page}, lu dans le .aux : la clé s'ajoute aux clés de la page.
\cs_new_protected:Npn \rlcite #1#2
  {
    \prop_get:NnNF \g__rl_page_prop {#2} \l__rl_tmp_tl { \tl_clear:N \l__rl_tmp_tl }
    \clist_if_in:NnF \l__rl_tmp_tl {#1}
      {
        \clist_put_right:Nn \l__rl_tmp_tl {#1}
        \prop_gput:NnV \g__rl_page_prop {#2} \l__rl_tmp_tl
      }
  }

%% Les tournures, lues dans <job>-tournures.aux (coordonnées en bp, l'origine en haut à gauche,
%% y vers le haut : négatif dans la page).
%%   \rlfond{page}{couleur}{x}{y}{largeur}{hauteur}   fond sous les mots, (x, y) coin bas gauche
%%   \rltrait{page}{x}{y}{largeur}                    trait plein sous les mots
%%   \rlpointille{page}{x}{y}{nombre}                 trait pointillé gris, un pas tous les 3 bp
%%   \rlnum{page}{couleur}{x}{y}{numéros}             numéros dans la marge, alignés à droite en x
%%   \rltournure{page}{n}{genre}{texte}{raison}       ligne de la bande ; genre f, c ou v
\cs_new_protected:Npn \rlfond #1#2#3#4#5#6
  {
    \__rl_ajouter:Nnn \g__rl_dessin_prop {#1}
      { \put (#3,#4) { \textcolor {#2} { \rule {#5\unitlength} {#6\unitlength} } } }
  }
\cs_new_protected:Npn \rltrait #1#2#3#4
  {
    \__rl_ajouter:Nnn \g__rl_dessin_prop {#1}
      { \put (#2,#3) { \textcolor {rltrait} { \rule {#4\unitlength} {0.6\unitlength} } } }
  }
\cs_new_protected:Npn \rlpointille #1#2#3#4
  {
    \__rl_ajouter:Nnn \g__rl_dessin_prop {#1}
      { \multiput (#2,#3) (3,0) {#4} { \textcolor {rlgris} { \rule {1.6\unitlength} {0.5\unitlength} } } }
  }
\cs_new_protected:Npn \rlnum #1#2#3#4#5
  {
    \__rl_ajouter:Nnn \g__rl_dessin_prop {#1}
      { \put (#3,#4) { \makebox (0,0) [r] { \normalfont \sffamily \bfseries \tiny \textcolor {#2} {#5} } } }
  }
\cs_new_protected:Npn \rltournure #1#2#3#4#5
  { \__rl_ajouter:Nnn \g__rl_tournure_prop {#1} { \__rl_ligne_tournure:nnnn {#2} {#3} {#4} {#5} } }

%% Le numéro de chaque clé, pris au passage de \bibcite dans le .aux : natbib écrit
%% {numéro}{année}{{auteurs}}{{...}}, LaTeX l'étiquette seule (un numéro, ou ABK+06 en style
%% alpha), beamer l'étiquette dans un \hyperlink vers sa page de bibliographie. Les \bibcite
%% viennent dans l'ordre de la bibliographie : leur rang ordonne les lignes de la bande.
\cs_new_eq:NN \__rl_bibcite_orig:nn \bibcite
\cs_gset_protected:Npn \bibcite #1#2
  {
    \__rl_bibcite_orig:nn {#1} {#2}
    \prop_if_in:NnF \g__rl_rang_prop {#1}
      {
        \int_gincr:N \g__rl_rang_int
        \prop_gput:Nne \g__rl_rang_prop {#1} { \int_use:N \g__rl_rang_int }
      }
    \cs_if_exist:NTF \NAT@parse
      { \prop_gput:Nno \g__rl_num_prop {#1} { \tl_head:n {#2} } }
      { \prop_gput:Nnn \g__rl_num_prop {#1} {#2} }
  }

%% Écrit #1 dans le .aux à la sortie de la page, avec la page absolue en dernier argument.
\cs_new_protected:Npn \__rl_ecrire_page:n #1
  {
    \protected@write \@auxout {}
      { #1 { \exp_not:N \the \exp_not:N \ReadonlyShipoutCounter } }
  }

%% Chaque clé citée écrit \rlcite{clé}{page absolue} à la sortie de sa page.
\cs_new_protected:Npn \__rl_note:n #1
  {
    \@bsphack
    \clist_map_inline:nn {#1}
      { \__rl_ecrire_page:n { \token_to_str:N \rlcite {##1} } }
    \@esphack
  }

%% Branchement, une fois tous les paquets en place (natbib choisit son \@citex au début du
%% document) : toutes les citations passent par \@citex, avec deux arguments optionnels sous
%% natbib, un seul avec le \cite de LaTeX ; la liste des clés est le dernier argument. Les clés
%% sont notées même indéfinies, dès le premier passage. (Dans une macro : le code d'un crochet ne
%% dédouble pas les #.) Le .aux reçoit aussi la hauteur de la page, pour relecture-tournures.py.
\cs_new_protected:Npn \__rl_brancher:
  {
    \cs_gset_eq:NN \__rl_citex_orig:w \@citex
    \cs_if_exist:NTF \NAT@parse
      {
        \cs_gset_protected:Npn \@citex [##1] [##2] ##3
          { \__rl_citex_orig:w [{##1}] [{##2}] {##3} \__rl_note:n {##3} }
      }
      {
        \cs_gset_protected:Npn \@citex [##1] ##2
          { \__rl_citex_orig:w [{##1}] {##2} \__rl_note:n {##2} }
      }
    \iow_now:Ne \@auxout { \token_to_str:N \rlhauteur { \dim_use:N \paperheight } }
  }
\AddToHook { begindocument / end } { \__rl_brancher: }

%% Les pages de la bibliographie, où relecture-tournures.py ne cherche rien : la page de son
%% premier paragraphe (le titre ou la première entrée) et celle de sa fin. Rien n'est écrit entre
%% deux paragraphes, où une note dans la liste verticale changerait les espaces.
\AddToHook { env / thebibliography / begin }
  { \AddToHookNext { para / begin } { \__rl_ecrire_page:n { \token_to_str:N \rlbiblio {debut} } } }
\AddToHook { env / thebibliography / end }
  { \@bsphack \__rl_ecrire_page:n { \token_to_str:N \rlbiblio {fin} } \@esphack }

%% La bande d'une page : les références, dans l'ordre de la bibliographie, puis les tournures.
\cs_new:Npn \__rl_num:n #1
  { \prop_if_in:NnTF \g__rl_rang_prop {#1} { \prop_item:Nn \g__rl_rang_prop {#1} } { 0 } }

%% L'étiquette entre crochets, lien vers la bibliographie ; celle de beamer porte déjà le sien
%% (deux liens emboîtés arrêteraient pdfTeX).
\cs_new_protected:Npn \__rl_etiquette:n #1
  {
    \prop_get:NnNF \g__rl_num_prop {#1} \l__rl_tmp_tl { \tl_set:Nn \l__rl_tmp_tl { ? } }
    \cs_if_exist:NTF \hyperlink
      {
        \tl_if_in:NnTF \l__rl_tmp_tl { \hyperlink }
          { [ \l__rl_tmp_tl ] }
          { \hyperlink { cite.#1 } { [ \l__rl_tmp_tl ] } }
      }
      { [ \l__rl_tmp_tl ] }
  }

\cs_new_protected:Npn \__rl_ligne:n #1
  {
    \hangindent = 2.9em \hangafter = 1 \noindent
    \hbox_set:Nn \l__rl_etiquette_box { \__rl_etiquette:n {#1} }
    \hbox_to_wd:nn { \dim_max:nn { 2.9em } { \box_wd:N \l__rl_etiquette_box + 0.5em } }
      { \box_use_drop:N \l__rl_etiquette_box \hfil }
    \prop_get:NnNTF \g__rl_entree_prop {#1} \l__rl_tmp_tl
      { \exp_after:wN \__rl_entree:nnn \l__rl_tmp_tl }
      { \texttt { \tl_to_str:n {#1} } }
    \par
  }
\cs_new_protected:Npn \__rl_entree:nnn #1#2#3
  { \tl_if_empty:nF {#1} { #1 \c_space_tl } \textit {#2} \tl_if_empty:nF {#3} { \c_space_tl #3 } }

%% Une ligne des tournures : numéro (celui de la marge), texte repéré, raison ; en gris pour un
%% indice faible.
\cs_new_protected:Npn \__rl_ligne_tournure:nnnn #1#2#3#4
  {
    \group_begin:
      \str_if_eq:nnT {#2} {v} { \color {rlgris} }
      \hangindent = 2.9em \hangafter = 1 \noindent
      \hbox_to_wd:nn { 2.9em }
        { \sffamily \bfseries \str_if_eq:nnF {#2} {v} { \color {rltrait} } #1 \hfil }
      \textit {#3} \c_space_tl : \c_space_tl #4
      \par
    \group_end:
  }

\cs_new_protected:Npn \__rl_bande:
  {
    \tl_if_empty:NF \l__rl_refs_tl
      {
        \exp_args:NNV \seq_set_from_clist:Nn \l__rl_cles_seq \l__rl_refs_tl
        \seq_sort:Nn \l__rl_cles_seq
          {
            \int_compare:nNnTF { \__rl_num:n {##1} } > { \__rl_num:n {##2} }
              { \sort_return_swapped: } { \sort_return_same: }
          }
      }
    \vbox_gset:Nn \g__rl_bande_box
      {
        \hsize = \textwidth \linewidth = \textwidth
        \normalfont \footnotesize \raggedright
        \parindent = 0pt \parskip = 1.5pt \relax
        \cs_set:Npn \\ { \relax \ignorespaces }
        \tl_if_empty:NF \l__rl_refs_tl
          {
            \noindent \textbf { \rltitre } \par
            \seq_map_function:NN \l__rl_cles_seq \__rl_ligne:n
          }
        \tl_if_empty:NF \l__rl_tour_tl
          {
            \tl_if_empty:NF \l__rl_refs_tl { \skip_vertical:n { 4pt } }
            \noindent \textbf { \rltitretournures } \par
            \l__rl_tour_tl
          }
      }
  }

%% Avant la sortie de chaque page (le compteur vaut alors le numéro de la page précédente) :
%% la bande, le dessin des tournures, et la hauteur de la page.
\AddToHook { shipout / before }
  {
    \dim_compare:nNnT { \g__rl_hauteur_dim } = { 0pt }
      { \dim_gset:Nn \g__rl_hauteur_dim { \paperheight } }
    \box_gclear:N \g__rl_bande_box
    \dim_gzero:N \g__rl_bande_dim
    \tl_set:Ne \l__rl_page_tl { \int_eval:n { \ReadonlyShipoutCounter + 1 } }
    \prop_get:NeNF \g__rl_page_prop { \l__rl_page_tl } \l__rl_refs_tl { \tl_clear:N \l__rl_refs_tl }
    \prop_get:NeNF \g__rl_tournure_prop { \l__rl_page_tl } \l__rl_tour_tl { \tl_clear:N \l__rl_tour_tl }
    \prop_get:NeNF \g__rl_dessin_prop { \l__rl_page_tl } \l__rl_tmp_tl { \tl_clear:N \l__rl_tmp_tl }
    \tl_gset_eq:NN \g__rl_dessin_tl \l__rl_tmp_tl
    \bool_lazy_and:nnF { \tl_if_empty_p:N \l__rl_refs_tl } { \tl_if_empty_p:N \l__rl_tour_tl }
      {
        \__rl_bande:
        \dim_gset:Nn \g__rl_bande_dim
          { \box_ht_plus_dp:N \g__rl_bande_box + 2 \c__rl_marge_dim }
      }
    \cs_if_exist:NT \pdfpageheight
      { \global \pdfpageheight = \dim_eval:n { \g__rl_hauteur_dim + \g__rl_bande_dim } \relax }
  }

%% Sous le texte de la page : les fonds et les traits des tournures, puis la bande, sur fond gris
%% clair et sur toute la largeur, son texte aligné sur la colonne.
\cs_new:Npn \__rl_pt:n #1 { \dim_to_decimal_in_unit:nn {#1} { 1pt } }
\AddToHook { shipout / background }
  {
    \tl_if_empty:NF \g__rl_dessin_tl
      {
        \group_begin:
          \setlength \unitlength { 1bp }
          \g__rl_dessin_tl
        \group_end:
      }
    \box_if_empty:NF \g__rl_bande_box
      {
        \setlength \unitlength { 1pt }
        \put ( 0 , \__rl_pt:n { - \g__rl_hauteur_dim - \g__rl_bande_dim } )
          { \textcolor { black!6 } { \rule { \paperwidth } { \g__rl_bande_dim } } }
        \put ( 0 , \__rl_pt:n { - \g__rl_hauteur_dim } )
          { \textcolor { black!35 } { \rule { \paperwidth } { 0.6pt } } }
        \put
          (
            \__rl_pt:n
              {
                1in + \hoffset +
                \bool_lazy_and:nnTF { \legacy_if_p:n { @twoside } } { \int_if_even_p:n { \c@page } }
                  { \evensidemargin } { \oddsidemargin }
              } ,
            \__rl_pt:n
              { - \g__rl_hauteur_dim - \c__rl_marge_dim - \box_ht:N \g__rl_bande_box }
          )
          { \box_use:N \g__rl_bande_box }
      }
  }

\makeatother
\ExplSyntaxOff

%% Les entrées courtes (absentes au premier passage) et, au dernier passage, les tournures, lues
%% hors de la syntaxe expl3 qui ignore les espaces. Comme pour la bibliographie, & et \\ sont
%% neutralisés (un champ journal en contient parfois, dans la biblio de référence).
\begingroup
  \catcode`\&=12
  \InputIfFileExists{\jobname-refs.bbl}{}{}
\endgroup
\ifdefined\rlavectournures
  \InputIfFileExists{\jobname-tournures.aux}{}{}
\fi
}\input{<source>}"
%% Références : chaque citation écrit \rlcite{clé}{page} dans le .aux, à la sortie de la page, donc
%% lu au passage suivant ; les entrées courtes viennent de <job>-refs.bbl, que bibtex produit avec
%% relecture.bst sur les clés citées. Marche avec natbib (acmart, LIPIcs) et avec le \cite de
%% LaTeX, hyperref chargé ou non ; pas avec biblatex.
%% Tournures : au dernier passage seulement (\rlavectournures défini), <job>-tournures.aux, écrit
%% par SCRIPTS/relecture-tournures.py d'après le texte du PDF (pdftotext -bbox) : des rectangles en
%% points PostScript (bp) depuis le coin haut gauche de la page, et les lignes de la bande. Pour
%% lui, le .aux donne la hauteur de la page sans bande (\rlhauteur) et les pages de la
%% bibliographie (\rlbiblio), où rien n'est cherché.
%% Ne part jamais vers le finalisé : le PDF de relecture est <document>-relecture.pdf.

\providecommand\rltitre{R\'ef\'erences cit\'ees sur cette page}
\providecommand\rltitretournures{Tournures \`a revoir}
\providecommand\rlhauteur[1]{}
\providecommand\rlbiblio[2]{}

\RequirePackage{xcolor}
\definecolor{rlfort}{RGB}{255,232,120}   % fond : tournure d'assistant (jaune)
\definecolor{rlfaute}{RGB}{196,222,255}  % fond : faute ou calque (bleu clair)
\definecolor{rltrait}{RGB}{25,55,145}    % traits pleins et numéros (bleu foncé)
\definecolor{rlgris}{gray}{0.45}         % indice faible : pointillé, numéro et ligne de la bande

\ExplSyntaxOn
\makeatletter

\prop_new:N \g__rl_entree_prop     % clé -> {Auteurs.}{Titre.}{Lieu, année.}
\prop_new:N \g__rl_num_prop        % clé -> étiquette dans la bibliographie
\prop_new:N \g__rl_rang_prop       % clé -> rang dans la bibliographie
\int_new:N \g__rl_rang_int
\box_new:N \l__rl_etiquette_box
\prop_new:N \g__rl_page_prop       % page absolue -> clés citées, séparées par des virgules
\prop_new:N \g__rl_dessin_prop     % page absolue -> fonds, traits et numéros des tournures
\prop_new:N \g__rl_tournure_prop   % page absolue -> lignes « Tournures à revoir » de la bande
\box_new:N \g__rl_bande_box
\tl_new:N \g__rl_dessin_tl         % dessin de la page en cours de sortie
\dim_new:N \g__rl_bande_dim        % hauteur ajoutée à la page courante
\dim_new:N \g__rl_hauteur_dim      % hauteur de la page sans bande
\seq_new:N \l__rl_cles_seq
\tl_new:N \l__rl_tmp_tl
\tl_new:N \l__rl_refs_tl
\tl_new:N \l__rl_tour_tl
\tl_new:N \l__rl_page_tl
\dim_const:Nn \c__rl_marge_dim { 14pt }
\cs_generate_variant:Nn \prop_get:NnNF { Ne }
\cs_generate_variant:Nn \prop_get:NnNT { Ne }

%% Ajoute du code à la valeur d'une clé (une page) d'une propriété globale.
\cs_new_protected:Npn \__rl_ajouter:Nnn #1#2#3
  {
    \prop_get:NnNF #1 {#2} \l__rl_tmp_tl { \tl_clear:N \l__rl_tmp_tl }
    \tl_put_right:Nn \l__rl_tmp_tl {#3}
    \prop_gput:NnV #1 {#2} \l__rl_tmp_tl
  }

%% Les entrées courtes, lues dans <job>-refs.bbl.
\cs_new_protected:Npn \rlentry #1#2#3#4
  { \prop_gput:Nnn \g__rl_entree_prop {#1} { {#2} {#3} {#4} } }

%% \rlcite{clé}{page}, lu dans le .aux : la clé s'ajoute aux clés de la page.
\cs_new_protected:Npn \rlcite #1#2
  {
    \prop_get:NnNF \g__rl_page_prop {#2} \l__rl_tmp_tl { \tl_clear:N \l__rl_tmp_tl }
    \clist_if_in:NnF \l__rl_tmp_tl {#1}
      {
        \clist_put_right:Nn \l__rl_tmp_tl {#1}
        \prop_gput:NnV \g__rl_page_prop {#2} \l__rl_tmp_tl
      }
  }

%% Les tournures, lues dans <job>-tournures.aux (coordonnées en bp, l'origine en haut à gauche,
%% y vers le haut : négatif dans la page).
%%   \rlfond{page}{couleur}{x}{y}{largeur}{hauteur}   fond sous les mots, (x, y) coin bas gauche
%%   \rltrait{page}{x}{y}{largeur}                    trait plein sous les mots
%%   \rlpointille{page}{x}{y}{nombre}                 trait pointillé gris, un pas tous les 3 bp
%%   \rlnum{page}{couleur}{x}{y}{numéros}             numéros dans la marge, alignés à droite en x
%%   \rltournure{page}{n}{genre}{texte}{raison}       ligne de la bande ; genre f, c ou v
\cs_new_protected:Npn \rlfond #1#2#3#4#5#6
  {
    \__rl_ajouter:Nnn \g__rl_dessin_prop {#1}
      { \put (#3,#4) { \textcolor {#2} { \rule {#5\unitlength} {#6\unitlength} } } }
  }
\cs_new_protected:Npn \rltrait #1#2#3#4
  {
    \__rl_ajouter:Nnn \g__rl_dessin_prop {#1}
      { \put (#2,#3) { \textcolor {rltrait} { \rule {#4\unitlength} {0.6\unitlength} } } }
  }
\cs_new_protected:Npn \rlpointille #1#2#3#4
  {
    \__rl_ajouter:Nnn \g__rl_dessin_prop {#1}
      { \multiput (#2,#3) (3,0) {#4} { \textcolor {rlgris} { \rule {1.6\unitlength} {0.5\unitlength} } } }
  }
\cs_new_protected:Npn \rlnum #1#2#3#4#5
  {
    \__rl_ajouter:Nnn \g__rl_dessin_prop {#1}
      { \put (#3,#4) { \makebox (0,0) [r] { \normalfont \sffamily \bfseries \tiny \textcolor {#2} {#5} } } }
  }
\cs_new_protected:Npn \rltournure #1#2#3#4#5
  { \__rl_ajouter:Nnn \g__rl_tournure_prop {#1} { \__rl_ligne_tournure:nnnn {#2} {#3} {#4} {#5} } }

%% Le numéro de chaque clé, pris au passage de \bibcite dans le .aux : natbib écrit
%% {numéro}{année}{{auteurs}}{{...}}, LaTeX l'étiquette seule (un numéro, ou ABK+06 en style
%% alpha), beamer l'étiquette dans un \hyperlink vers sa page de bibliographie. Les \bibcite
%% viennent dans l'ordre de la bibliographie : leur rang ordonne les lignes de la bande.
\cs_new_eq:NN \__rl_bibcite_orig:nn \bibcite
\cs_gset_protected:Npn \bibcite #1#2
  {
    \__rl_bibcite_orig:nn {#1} {#2}
    \prop_if_in:NnF \g__rl_rang_prop {#1}
      {
        \int_gincr:N \g__rl_rang_int
        \prop_gput:Nne \g__rl_rang_prop {#1} { \int_use:N \g__rl_rang_int }
      }
    \cs_if_exist:NTF \NAT@parse
      { \prop_gput:Nno \g__rl_num_prop {#1} { \tl_head:n {#2} } }
      { \prop_gput:Nnn \g__rl_num_prop {#1} {#2} }
  }

%% Écrit #1 dans le .aux à la sortie de la page, avec la page absolue en dernier argument.
\cs_new_protected:Npn \__rl_ecrire_page:n #1
  {
    \protected@write \@auxout {}
      { #1 { \exp_not:N \the \exp_not:N \ReadonlyShipoutCounter } }
  }

%% Chaque clé citée écrit \rlcite{clé}{page absolue} à la sortie de sa page.
\cs_new_protected:Npn \__rl_note:n #1
  {
    \@bsphack
    \clist_map_inline:nn {#1}
      { \__rl_ecrire_page:n { \token_to_str:N \rlcite {##1} } }
    \@esphack
  }

%% Branchement, une fois tous les paquets en place (natbib choisit son \@citex au début du
%% document) : toutes les citations passent par \@citex, avec deux arguments optionnels sous
%% natbib, un seul avec le \cite de LaTeX ; la liste des clés est le dernier argument. Les clés
%% sont notées même indéfinies, dès le premier passage. (Dans une macro : le code d'un crochet ne
%% dédouble pas les #.) Le .aux reçoit aussi la hauteur de la page, pour relecture-tournures.py.
\cs_new_protected:Npn \__rl_brancher:
  {
    \cs_gset_eq:NN \__rl_citex_orig:w \@citex
    \cs_if_exist:NTF \NAT@parse
      {
        \cs_gset_protected:Npn \@citex [##1] [##2] ##3
          { \__rl_citex_orig:w [{##1}] [{##2}] {##3} \__rl_note:n {##3} }
      }
      {
        \cs_gset_protected:Npn \@citex [##1] ##2
          { \__rl_citex_orig:w [{##1}] {##2} \__rl_note:n {##2} }
      }
    \iow_now:Ne \@auxout { \token_to_str:N \rlhauteur { \dim_use:N \paperheight } }
  }
\AddToHook { begindocument / end } { \__rl_brancher: }

%% Les pages de la bibliographie, où relecture-tournures.py ne cherche rien : la page de son
%% premier paragraphe (le titre ou la première entrée) et celle de sa fin. Rien n'est écrit entre
%% deux paragraphes, où une note dans la liste verticale changerait les espaces.
\AddToHook { env / thebibliography / begin }
  { \AddToHookNext { para / begin } { \__rl_ecrire_page:n { \token_to_str:N \rlbiblio {debut} } } }
\AddToHook { env / thebibliography / end }
  { \@bsphack \__rl_ecrire_page:n { \token_to_str:N \rlbiblio {fin} } \@esphack }

%% La bande d'une page : les références, dans l'ordre de la bibliographie, puis les tournures.
\cs_new:Npn \__rl_num:n #1
  { \prop_if_in:NnTF \g__rl_rang_prop {#1} { \prop_item:Nn \g__rl_rang_prop {#1} } { 0 } }

%% L'étiquette entre crochets, lien vers la bibliographie ; celle de beamer porte déjà le sien
%% (deux liens emboîtés arrêteraient pdfTeX).
\cs_new_protected:Npn \__rl_etiquette:n #1
  {
    \prop_get:NnNF \g__rl_num_prop {#1} \l__rl_tmp_tl { \tl_set:Nn \l__rl_tmp_tl { ? } }
    \cs_if_exist:NTF \hyperlink
      {
        \tl_if_in:NnTF \l__rl_tmp_tl { \hyperlink }
          { [ \l__rl_tmp_tl ] }
          { \hyperlink { cite.#1 } { [ \l__rl_tmp_tl ] } }
      }
      { [ \l__rl_tmp_tl ] }
  }

\cs_new_protected:Npn \__rl_ligne:n #1
  {
    \hangindent = 2.9em \hangafter = 1 \noindent
    \hbox_set:Nn \l__rl_etiquette_box { \__rl_etiquette:n {#1} }
    \hbox_to_wd:nn { \dim_max:nn { 2.9em } { \box_wd:N \l__rl_etiquette_box + 0.5em } }
      { \box_use_drop:N \l__rl_etiquette_box \hfil }
    \prop_get:NnNTF \g__rl_entree_prop {#1} \l__rl_tmp_tl
      { \exp_after:wN \__rl_entree:nnn \l__rl_tmp_tl }
      { \texttt { \tl_to_str:n {#1} } }
    \par
  }
\cs_new_protected:Npn \__rl_entree:nnn #1#2#3
  { \tl_if_empty:nF {#1} { #1 \c_space_tl } \textit {#2} \tl_if_empty:nF {#3} { \c_space_tl #3 } }

%% Une ligne des tournures : numéro (celui de la marge), texte repéré, raison ; en gris pour un
%% indice faible.
\cs_new_protected:Npn \__rl_ligne_tournure:nnnn #1#2#3#4
  {
    \group_begin:
      \str_if_eq:nnT {#2} {v} { \color {rlgris} }
      \hangindent = 2.9em \hangafter = 1 \noindent
      \hbox_to_wd:nn { 2.9em }
        { \sffamily \bfseries \str_if_eq:nnF {#2} {v} { \color {rltrait} } #1 \hfil }
      \textit {#3} \c_space_tl : \c_space_tl #4
      \par
    \group_end:
  }

\cs_new_protected:Npn \__rl_bande:
  {
    \tl_if_empty:NF \l__rl_refs_tl
      {
        \exp_args:NNV \seq_set_from_clist:Nn \l__rl_cles_seq \l__rl_refs_tl
        \seq_sort:Nn \l__rl_cles_seq
          {
            \int_compare:nNnTF { \__rl_num:n {##1} } > { \__rl_num:n {##2} }
              { \sort_return_swapped: } { \sort_return_same: }
          }
      }
    \vbox_gset:Nn \g__rl_bande_box
      {
        \hsize = \textwidth \linewidth = \textwidth
        \normalfont \footnotesize \raggedright
        \parindent = 0pt \parskip = 1.5pt \relax
        \cs_set:Npn \\ { \relax \ignorespaces }
        \tl_if_empty:NF \l__rl_refs_tl
          {
            \noindent \textbf { \rltitre } \par
            \seq_map_function:NN \l__rl_cles_seq \__rl_ligne:n
          }
        \tl_if_empty:NF \l__rl_tour_tl
          {
            \tl_if_empty:NF \l__rl_refs_tl { \skip_vertical:n { 4pt } }
            \noindent \textbf { \rltitretournures } \par
            \l__rl_tour_tl
          }
      }
  }

%% Avant la sortie de chaque page (le compteur vaut alors le numéro de la page précédente) :
%% la bande, le dessin des tournures, et la hauteur de la page.
\AddToHook { shipout / before }
  {
    \dim_compare:nNnT { \g__rl_hauteur_dim } = { 0pt }
      { \dim_gset:Nn \g__rl_hauteur_dim { \paperheight } }
    \box_gclear:N \g__rl_bande_box
    \dim_gzero:N \g__rl_bande_dim
    \tl_set:Ne \l__rl_page_tl { \int_eval:n { \ReadonlyShipoutCounter + 1 } }
    \prop_get:NeNF \g__rl_page_prop { \l__rl_page_tl } \l__rl_refs_tl { \tl_clear:N \l__rl_refs_tl }
    \prop_get:NeNF \g__rl_tournure_prop { \l__rl_page_tl } \l__rl_tour_tl { \tl_clear:N \l__rl_tour_tl }
    \prop_get:NeNF \g__rl_dessin_prop { \l__rl_page_tl } \l__rl_tmp_tl { \tl_clear:N \l__rl_tmp_tl }
    \tl_gset_eq:NN \g__rl_dessin_tl \l__rl_tmp_tl
    \bool_lazy_and:nnF { \tl_if_empty_p:N \l__rl_refs_tl } { \tl_if_empty_p:N \l__rl_tour_tl }
      {
        \__rl_bande:
        \dim_gset:Nn \g__rl_bande_dim
          { \box_ht_plus_dp:N \g__rl_bande_box + 2 \c__rl_marge_dim }
      }
    \cs_if_exist:NT \pdfpageheight
      { \global \pdfpageheight = \dim_eval:n { \g__rl_hauteur_dim + \g__rl_bande_dim } \relax }
  }

%% Sous le texte de la page : les fonds et les traits des tournures, puis la bande, sur fond gris
%% clair et sur toute la largeur, son texte aligné sur la colonne.
\cs_new:Npn \__rl_pt:n #1 { \dim_to_decimal_in_unit:nn {#1} { 1pt } }
\AddToHook { shipout / background }
  {
    \tl_if_empty:NF \g__rl_dessin_tl
      {
        \group_begin:
          \setlength \unitlength { 1bp }
          \g__rl_dessin_tl
        \group_end:
      }
    \box_if_empty:NF \g__rl_bande_box
      {
        \setlength \unitlength { 1pt }
        \put ( 0 , \__rl_pt:n { - \g__rl_hauteur_dim - \g__rl_bande_dim } )
          { \textcolor { black!6 } { \rule { \paperwidth } { \g__rl_bande_dim } } }
        \put ( 0 , \__rl_pt:n { - \g__rl_hauteur_dim } )
          { \textcolor { black!35 } { \rule { \paperwidth } { 0.6pt } } }
        \put
          (
            \__rl_pt:n
              {
                1in + \hoffset +
                \bool_lazy_and:nnTF { \legacy_if_p:n { @twoside } } { \int_if_even_p:n { \c@page } }
                  { \evensidemargin } { \oddsidemargin }
              } ,
            \__rl_pt:n
              { - \g__rl_hauteur_dim - \c__rl_marge_dim - \box_ht:N \g__rl_bande_box }
          )
          { \box_use:N \g__rl_bande_box }
      }
  }

\makeatother
\ExplSyntaxOff

%% Les entrées courtes (absentes au premier passage) et, au dernier passage, les tournures, lues
%% hors de la syntaxe expl3 qui ignore les espaces. Comme pour la bibliographie, & et \\ sont
%% neutralisés (un champ journal en contient parfois, dans la biblio de référence).
\begingroup
  \catcode`\&=12
  \InputIfFileExists{\jobname-refs.bbl}{}{}
\endgroup
\ifdefined\rlavectournures
  \InputIfFileExists{\jobname-tournures.aux}{}{}
\fi
}\input{<source>}"
%% Références : chaque citation écrit \rlcite{clé}{page} dans le .aux, à la sortie de la page, donc
%% lu au passage suivant ; les entrées courtes viennent de <job>-refs.bbl, que bibtex produit avec
%% relecture.bst sur les clés citées. Marche avec natbib (acmart, LIPIcs) et avec le \cite de
%% LaTeX, hyperref chargé ou non ; pas avec biblatex.
%% Tournures : au dernier passage seulement (\rlavectournures défini), <job>-tournures.aux, écrit
%% par SCRIPTS/relecture-tournures.py d'après le texte du PDF (pdftotext -bbox) : des rectangles en
%% points PostScript (bp) depuis le coin haut gauche de la page, et les lignes de la bande. Pour
%% lui, le .aux donne la hauteur de la page sans bande (\rlhauteur) et les pages de la
%% bibliographie (\rlbiblio), où rien n'est cherché.
%% Ne part jamais vers le finalisé : le PDF de relecture est <document>-relecture.pdf.

\providecommand\rltitre{R\'ef\'erences cit\'ees sur cette page}
\providecommand\rltitretournures{Tournures \`a revoir}
\providecommand\rlhauteur[1]{}
\providecommand\rlbiblio[2]{}

\RequirePackage{xcolor}
\definecolor{rlfort}{RGB}{255,232,120}   % fond : tournure d'assistant (jaune)
\definecolor{rlfaute}{RGB}{196,222,255}  % fond : faute ou calque (bleu clair)
\definecolor{rltrait}{RGB}{25,55,145}    % traits pleins et numéros (bleu foncé)
\definecolor{rlgris}{gray}{0.45}         % indice faible : pointillé, numéro et ligne de la bande

\ExplSyntaxOn
\makeatletter

\prop_new:N \g__rl_entree_prop     % clé -> {Auteurs.}{Titre.}{Lieu, année.}
\prop_new:N \g__rl_num_prop        % clé -> étiquette dans la bibliographie
\prop_new:N \g__rl_rang_prop       % clé -> rang dans la bibliographie
\int_new:N \g__rl_rang_int
\box_new:N \l__rl_etiquette_box
\prop_new:N \g__rl_page_prop       % page absolue -> clés citées, séparées par des virgules
\prop_new:N \g__rl_dessin_prop     % page absolue -> fonds, traits et numéros des tournures
\prop_new:N \g__rl_tournure_prop   % page absolue -> lignes « Tournures à revoir » de la bande
\box_new:N \g__rl_bande_box
\tl_new:N \g__rl_dessin_tl         % dessin de la page en cours de sortie
\dim_new:N \g__rl_bande_dim        % hauteur ajoutée à la page courante
\dim_new:N \g__rl_hauteur_dim      % hauteur de la page sans bande
\seq_new:N \l__rl_cles_seq
\tl_new:N \l__rl_tmp_tl
\tl_new:N \l__rl_refs_tl
\tl_new:N \l__rl_tour_tl
\tl_new:N \l__rl_page_tl
\dim_const:Nn \c__rl_marge_dim { 14pt }
\cs_generate_variant:Nn \prop_get:NnNF { Ne }
\cs_generate_variant:Nn \prop_get:NnNT { Ne }

%% Ajoute du code à la valeur d'une clé (une page) d'une propriété globale.
\cs_new_protected:Npn \__rl_ajouter:Nnn #1#2#3
  {
    \prop_get:NnNF #1 {#2} \l__rl_tmp_tl { \tl_clear:N \l__rl_tmp_tl }
    \tl_put_right:Nn \l__rl_tmp_tl {#3}
    \prop_gput:NnV #1 {#2} \l__rl_tmp_tl
  }

%% Les entrées courtes, lues dans <job>-refs.bbl.
\cs_new_protected:Npn \rlentry #1#2#3#4
  { \prop_gput:Nnn \g__rl_entree_prop {#1} { {#2} {#3} {#4} } }

%% \rlcite{clé}{page}, lu dans le .aux : la clé s'ajoute aux clés de la page.
\cs_new_protected:Npn \rlcite #1#2
  {
    \prop_get:NnNF \g__rl_page_prop {#2} \l__rl_tmp_tl { \tl_clear:N \l__rl_tmp_tl }
    \clist_if_in:NnF \l__rl_tmp_tl {#1}
      {
        \clist_put_right:Nn \l__rl_tmp_tl {#1}
        \prop_gput:NnV \g__rl_page_prop {#2} \l__rl_tmp_tl
      }
  }

%% Les tournures, lues dans <job>-tournures.aux (coordonnées en bp, l'origine en haut à gauche,
%% y vers le haut : négatif dans la page).
%%   \rlfond{page}{couleur}{x}{y}{largeur}{hauteur}   fond sous les mots, (x, y) coin bas gauche
%%   \rltrait{page}{x}{y}{largeur}                    trait plein sous les mots
%%   \rlpointille{page}{x}{y}{nombre}                 trait pointillé gris, un pas tous les 3 bp
%%   \rlnum{page}{couleur}{x}{y}{numéros}             numéros dans la marge, alignés à droite en x
%%   \rltournure{page}{n}{genre}{texte}{raison}       ligne de la bande ; genre f, c ou v
\cs_new_protected:Npn \rlfond #1#2#3#4#5#6
  {
    \__rl_ajouter:Nnn \g__rl_dessin_prop {#1}
      { \put (#3,#4) { \textcolor {#2} { \rule {#5\unitlength} {#6\unitlength} } } }
  }
\cs_new_protected:Npn \rltrait #1#2#3#4
  {
    \__rl_ajouter:Nnn \g__rl_dessin_prop {#1}
      { \put (#2,#3) { \textcolor {rltrait} { \rule {#4\unitlength} {0.6\unitlength} } } }
  }
\cs_new_protected:Npn \rlpointille #1#2#3#4
  {
    \__rl_ajouter:Nnn \g__rl_dessin_prop {#1}
      { \multiput (#2,#3) (3,0) {#4} { \textcolor {rlgris} { \rule {1.6\unitlength} {0.5\unitlength} } } }
  }
\cs_new_protected:Npn \rlnum #1#2#3#4#5
  {
    \__rl_ajouter:Nnn \g__rl_dessin_prop {#1}
      { \put (#3,#4) { \makebox (0,0) [r] { \normalfont \sffamily \bfseries \tiny \textcolor {#2} {#5} } } }
  }
\cs_new_protected:Npn \rltournure #1#2#3#4#5
  { \__rl_ajouter:Nnn \g__rl_tournure_prop {#1} { \__rl_ligne_tournure:nnnn {#2} {#3} {#4} {#5} } }

%% Le numéro de chaque clé, pris au passage de \bibcite dans le .aux : natbib écrit
%% {numéro}{année}{{auteurs}}{{...}}, LaTeX l'étiquette seule (un numéro, ou ABK+06 en style
%% alpha), beamer l'étiquette dans un \hyperlink vers sa page de bibliographie. Les \bibcite
%% viennent dans l'ordre de la bibliographie : leur rang ordonne les lignes de la bande.
\cs_new_eq:NN \__rl_bibcite_orig:nn \bibcite
\cs_gset_protected:Npn \bibcite #1#2
  {
    \__rl_bibcite_orig:nn {#1} {#2}
    \prop_if_in:NnF \g__rl_rang_prop {#1}
      {
        \int_gincr:N \g__rl_rang_int
        \prop_gput:Nne \g__rl_rang_prop {#1} { \int_use:N \g__rl_rang_int }
      }
    \cs_if_exist:NTF \NAT@parse
      { \prop_gput:Nno \g__rl_num_prop {#1} { \tl_head:n {#2} } }
      { \prop_gput:Nnn \g__rl_num_prop {#1} {#2} }
  }

%% Écrit #1 dans le .aux à la sortie de la page, avec la page absolue en dernier argument.
\cs_new_protected:Npn \__rl_ecrire_page:n #1
  {
    \protected@write \@auxout {}
      { #1 { \exp_not:N \the \exp_not:N \ReadonlyShipoutCounter } }
  }

%% Chaque clé citée écrit \rlcite{clé}{page absolue} à la sortie de sa page.
\cs_new_protected:Npn \__rl_note:n #1
  {
    \@bsphack
    \clist_map_inline:nn {#1}
      { \__rl_ecrire_page:n { \token_to_str:N \rlcite {##1} } }
    \@esphack
  }

%% Branchement, une fois tous les paquets en place (natbib choisit son \@citex au début du
%% document) : toutes les citations passent par \@citex, avec deux arguments optionnels sous
%% natbib, un seul avec le \cite de LaTeX ; la liste des clés est le dernier argument. Les clés
%% sont notées même indéfinies, dès le premier passage. (Dans une macro : le code d'un crochet ne
%% dédouble pas les #.) Le .aux reçoit aussi la hauteur de la page, pour relecture-tournures.py.
\cs_new_protected:Npn \__rl_brancher:
  {
    \cs_gset_eq:NN \__rl_citex_orig:w \@citex
    \cs_if_exist:NTF \NAT@parse
      {
        \cs_gset_protected:Npn \@citex [##1] [##2] ##3
          { \__rl_citex_orig:w [{##1}] [{##2}] {##3} \__rl_note:n {##3} }
      }
      {
        \cs_gset_protected:Npn \@citex [##1] ##2
          { \__rl_citex_orig:w [{##1}] {##2} \__rl_note:n {##2} }
      }
    \iow_now:Ne \@auxout { \token_to_str:N \rlhauteur { \dim_use:N \paperheight } }
  }
\AddToHook { begindocument / end } { \__rl_brancher: }

%% Les pages de la bibliographie, où relecture-tournures.py ne cherche rien : la page de son
%% premier paragraphe (le titre ou la première entrée) et celle de sa fin. Rien n'est écrit entre
%% deux paragraphes, où une note dans la liste verticale changerait les espaces.
\AddToHook { env / thebibliography / begin }
  { \AddToHookNext { para / begin } { \__rl_ecrire_page:n { \token_to_str:N \rlbiblio {debut} } } }
\AddToHook { env / thebibliography / end }
  { \@bsphack \__rl_ecrire_page:n { \token_to_str:N \rlbiblio {fin} } \@esphack }

%% La bande d'une page : les références, dans l'ordre de la bibliographie, puis les tournures.
\cs_new:Npn \__rl_num:n #1
  { \prop_if_in:NnTF \g__rl_rang_prop {#1} { \prop_item:Nn \g__rl_rang_prop {#1} } { 0 } }

%% L'étiquette entre crochets, lien vers la bibliographie ; celle de beamer porte déjà le sien
%% (deux liens emboîtés arrêteraient pdfTeX).
\cs_new_protected:Npn \__rl_etiquette:n #1
  {
    \prop_get:NnNF \g__rl_num_prop {#1} \l__rl_tmp_tl { \tl_set:Nn \l__rl_tmp_tl { ? } }
    \cs_if_exist:NTF \hyperlink
      {
        \tl_if_in:NnTF \l__rl_tmp_tl { \hyperlink }
          { [ \l__rl_tmp_tl ] }
          { \hyperlink { cite.#1 } { [ \l__rl_tmp_tl ] } }
      }
      { [ \l__rl_tmp_tl ] }
  }

\cs_new_protected:Npn \__rl_ligne:n #1
  {
    \hangindent = 2.9em \hangafter = 1 \noindent
    \hbox_set:Nn \l__rl_etiquette_box { \__rl_etiquette:n {#1} }
    \hbox_to_wd:nn { \dim_max:nn { 2.9em } { \box_wd:N \l__rl_etiquette_box + 0.5em } }
      { \box_use_drop:N \l__rl_etiquette_box \hfil }
    \prop_get:NnNTF \g__rl_entree_prop {#1} \l__rl_tmp_tl
      { \exp_after:wN \__rl_entree:nnn \l__rl_tmp_tl }
      { \texttt { \tl_to_str:n {#1} } }
    \par
  }
\cs_new_protected:Npn \__rl_entree:nnn #1#2#3
  { \tl_if_empty:nF {#1} { #1 \c_space_tl } \textit {#2} \tl_if_empty:nF {#3} { \c_space_tl #3 } }

%% Une ligne des tournures : numéro (celui de la marge), texte repéré, raison ; en gris pour un
%% indice faible.
\cs_new_protected:Npn \__rl_ligne_tournure:nnnn #1#2#3#4
  {
    \group_begin:
      \str_if_eq:nnT {#2} {v} { \color {rlgris} }
      \hangindent = 2.9em \hangafter = 1 \noindent
      \hbox_to_wd:nn { 2.9em }
        { \sffamily \bfseries \str_if_eq:nnF {#2} {v} { \color {rltrait} } #1 \hfil }
      \textit {#3} \c_space_tl : \c_space_tl #4
      \par
    \group_end:
  }

\cs_new_protected:Npn \__rl_bande:
  {
    \tl_if_empty:NF \l__rl_refs_tl
      {
        \exp_args:NNV \seq_set_from_clist:Nn \l__rl_cles_seq \l__rl_refs_tl
        \seq_sort:Nn \l__rl_cles_seq
          {
            \int_compare:nNnTF { \__rl_num:n {##1} } > { \__rl_num:n {##2} }
              { \sort_return_swapped: } { \sort_return_same: }
          }
      }
    \vbox_gset:Nn \g__rl_bande_box
      {
        \hsize = \textwidth \linewidth = \textwidth
        \normalfont \footnotesize \raggedright
        \parindent = 0pt \parskip = 1.5pt \relax
        \cs_set:Npn \\ { \relax \ignorespaces }
        \tl_if_empty:NF \l__rl_refs_tl
          {
            \noindent \textbf { \rltitre } \par
            \seq_map_function:NN \l__rl_cles_seq \__rl_ligne:n
          }
        \tl_if_empty:NF \l__rl_tour_tl
          {
            \tl_if_empty:NF \l__rl_refs_tl { \skip_vertical:n { 4pt } }
            \noindent \textbf { \rltitretournures } \par
            \l__rl_tour_tl
          }
      }
  }

%% Avant la sortie de chaque page (le compteur vaut alors le numéro de la page précédente) :
%% la bande, le dessin des tournures, et la hauteur de la page.
\AddToHook { shipout / before }
  {
    \dim_compare:nNnT { \g__rl_hauteur_dim } = { 0pt }
      { \dim_gset:Nn \g__rl_hauteur_dim { \paperheight } }
    \box_gclear:N \g__rl_bande_box
    \dim_gzero:N \g__rl_bande_dim
    \tl_set:Ne \l__rl_page_tl { \int_eval:n { \ReadonlyShipoutCounter + 1 } }
    \prop_get:NeNF \g__rl_page_prop { \l__rl_page_tl } \l__rl_refs_tl { \tl_clear:N \l__rl_refs_tl }
    \prop_get:NeNF \g__rl_tournure_prop { \l__rl_page_tl } \l__rl_tour_tl { \tl_clear:N \l__rl_tour_tl }
    \prop_get:NeNF \g__rl_dessin_prop { \l__rl_page_tl } \l__rl_tmp_tl { \tl_clear:N \l__rl_tmp_tl }
    \tl_gset_eq:NN \g__rl_dessin_tl \l__rl_tmp_tl
    \bool_lazy_and:nnF { \tl_if_empty_p:N \l__rl_refs_tl } { \tl_if_empty_p:N \l__rl_tour_tl }
      {
        \__rl_bande:
        \dim_gset:Nn \g__rl_bande_dim
          { \box_ht_plus_dp:N \g__rl_bande_box + 2 \c__rl_marge_dim }
      }
    \cs_if_exist:NT \pdfpageheight
      { \global \pdfpageheight = \dim_eval:n { \g__rl_hauteur_dim + \g__rl_bande_dim } \relax }
  }

%% Sous le texte de la page : les fonds et les traits des tournures, puis la bande, sur fond gris
%% clair et sur toute la largeur, son texte aligné sur la colonne.
\cs_new:Npn \__rl_pt:n #1 { \dim_to_decimal_in_unit:nn {#1} { 1pt } }
\AddToHook { shipout / background }
  {
    \tl_if_empty:NF \g__rl_dessin_tl
      {
        \group_begin:
          \setlength \unitlength { 1bp }
          \g__rl_dessin_tl
        \group_end:
      }
    \box_if_empty:NF \g__rl_bande_box
      {
        \setlength \unitlength { 1pt }
        \put ( 0 , \__rl_pt:n { - \g__rl_hauteur_dim - \g__rl_bande_dim } )
          { \textcolor { black!6 } { \rule { \paperwidth } { \g__rl_bande_dim } } }
        \put ( 0 , \__rl_pt:n { - \g__rl_hauteur_dim } )
          { \textcolor { black!35 } { \rule { \paperwidth } { 0.6pt } } }
        \put
          (
            \__rl_pt:n
              {
                1in + \hoffset +
                \bool_lazy_and:nnTF { \legacy_if_p:n { @twoside } } { \int_if_even_p:n { \c@page } }
                  { \evensidemargin } { \oddsidemargin }
              } ,
            \__rl_pt:n
              { - \g__rl_hauteur_dim - \c__rl_marge_dim - \box_ht:N \g__rl_bande_box }
          )
          { \box_use:N \g__rl_bande_box }
      }
  }

\makeatother
\ExplSyntaxOff

%% Les entrées courtes (absentes au premier passage) et, au dernier passage, les tournures, lues
%% hors de la syntaxe expl3 qui ignore les espaces. Comme pour la bibliographie, & et \\ sont
%% neutralisés (un champ journal en contient parfois, dans la biblio de référence).
\begingroup
  \catcode`\&=12
  \InputIfFileExists{\jobname-refs.bbl}{}{}
\endgroup
\ifdefined\rlavectournures
  \InputIfFileExists{\jobname-tournures.aux}{}{}
\fi
